% ============================================================
\chapter{Version Control with Git and GitHub}
\label{ch:git}
% ============================================================

% ============================================================
\section{Configuring SSH Authentication for GitHub on Linux}
% ============================================================
GitHub no longer supports password authentication for Git operations over HTTPS. 
Instead, authentication must be performed using either a Personal Access Token or 
an SSH key. The recommended method for Linux environments is SSH authentication.

The following procedure generates an SSH key and configures Git to authenticate
with GitHub.

\subsection{Generate an SSH Key}

Create a new SSH key using the \texttt{ed25519} algorithm:

\begin{lstlisting}[basicstyle=\footnotesize\ttfamily,escapeinside={(*@}{@*)}]
ssh-keygen -t ed25519 -C "your_email@example.com"
\end{lstlisting}

When prompted for the file location, press \texttt{Enter} to accept the default:

\begin{lstlisting}[basicstyle=\footnotesize\ttfamily,escapeinside={(*@}{@*)}]
~/.ssh/id_ed25519
\end{lstlisting}

Optionally provide a passphrase for additional security.

\subsection{Start the SSH Agent}

Start the SSH agent and add the newly created key:

\begin{lstlisting}[basicstyle=\footnotesize\ttfamily,escapeinside={(*@}{@*)}]
eval "$(ssh-agent -s)"
ssh-add ~/.ssh/id_ed25519
\end{lstlisting}

\subsection{Add the Public Key to GitHub}

Display the public key:

\begin{lstlisting}[basicstyle=\footnotesize\ttfamily,escapeinside={(*@}{@*)}]
cat ~/.ssh/id_ed25519.pub
\end{lstlisting}

Copy the entire output and add it to GitHub:

\begin{enumerate}
\item Navigate to \texttt{GitHub → Settings → SSH and GPG keys}
\item Click \texttt{New SSH key}
\item Paste the copied key and save
\end{enumerate}

\subsection{Verify the Connection}

Test the authentication:

\begin{lstlisting}[basicstyle=\footnotesize\ttfamily,escapeinside={(*@}{@*)}]
ssh -T git@github.com
\end{lstlisting}

A successful configuration will display a message similar to:

\begin{lstlisting}[basicstyle=\footnotesize\ttfamily,escapeinside={(*@}{@*)}]
Hi <username>! You've successfully authenticated.
\end{lstlisting}

\subsection{Clone a Repository Using SSH}

Repositories must be cloned using the SSH URL:

\begin{lstlisting}[basicstyle=\footnotesize\ttfamily,escapeinside={(*@}{@*)}]
git clone git@github.com:USER/REPOSITORY.git
\end{lstlisting}

Using SSH avoids credential prompts and enables secure automated
Git operations.

% ============================================================
\section{Configuring SSH Authentication for GitHub on Windows}
\label{sec:ssh-windows}
% ============================================================

Git supports two protocols for communicating with remote repositories: HTTPS and SSH. While HTTPS is simpler to set up initially, SSH authentication offers a more seamless workflow for frequent contributors because it eliminates repeated credential prompts. This section walks through the complete process of configuring SSH-based authentication for GitHub on a Windows machine using Git Bash.

\subsection{Opening Git Bash}

Git for Windows ships with its own bundled OpenSSH client and a Unix-like shell called Git Bash. All SSH commands in this section must be run inside Git Bash, not in Command Prompt or PowerShell, because the bundled \texttt{ssh}, \texttt{ssh-keygen}, and \texttt{ssh-agent} executables are located in the Git for Windows installation directory and are not automatically available in other terminals.

To open Git Bash, right-click on the Desktop or any folder in File Explorer and select \textbf{Git Bash Here} from the context menu. Alternatively, search for \textbf{Git Bash} in the Windows Start menu and launch it from there.

\subsection{Generating an SSH Key}

SSH authentication relies on a public/private key pair. The private key remains on your machine; the public key is uploaded to GitHub. Generate a new key pair using the \texttt{ed25519} algorithm, which is modern, compact, and recommended over the older RSA default:

\begin{lstlisting}[style=bashstyle, caption={Generate an ed25519 SSH key pair},escapeinside={(*@}{@*)}]
ssh-keygen -t ed25519 -C "your_email@example.com"
\end{lstlisting}

Replace \texttt{your\_email@example.com} with the address associated with your GitHub account. When prompted for a file location, press \textbf{Enter} to accept the default path:

\begin{lstlisting}[style=windowsStyle, caption={Default key location on Windows},escapeinside={(*@}{@*)}]
C:\Users\YourName\.ssh\id_ed25519
\end{lstlisting}

You will then be asked to enter an optional passphrase. A passphrase adds a second layer of protection in case your private key file is ever exposed; however, it must be entered each time the key is used (unless you configure the SSH agent to cache it). Press \textbf{Enter} twice to skip the passphrase if you prefer not to set one.

\subsection{Starting the SSH Agent}

The SSH agent is a background process that holds decrypted private keys in memory, so you do not have to re-enter a passphrase for every Git operation. On Linux and macOS the agent is typically started automatically by the desktop session, but on Windows the \texttt{ssh-agent} bundled with Git Bash must be started manually in each new terminal session.

Start the agent and add your private key with the following two commands:

\begin{lstlisting}[style=bashstyle, caption={Start the SSH agent and add the private key},escapeinside={(*@}{@*)}]
eval "$(ssh-agent -s)"
ssh-add ~/.ssh/id_ed25519
\end{lstlisting}

The first command launches the agent and sets the necessary environment variables in the current shell. The second command registers your private key with the running agent. If you set a passphrase during key generation, you will be prompted to enter it now.

\begin{tipbox}{Automating Agent Startup}
The commands above must be repeated every time you open a new Git Bash session. To avoid this, you can add the following lines to your \texttt{\textasciitilde/.bashrc} file so that the agent starts automatically:

\begin{lstlisting}[style=bashstyle,escapeinside={(*@}{@*)}]
eval "$(ssh-agent -s)" > /dev/null 2>&1
ssh-add ~/.ssh/id_ed25519 2>/dev/null
\end{lstlisting}

Open or create \texttt{\textasciitilde/.bashrc} in a text editor, paste the two lines at the end of the file, and save. The redirects suppress startup messages so they do not clutter your terminal. The next time you open Git Bash the agent will start and load your key without any manual steps.
\end{tipbox}

\subsection{Adding the Public Key to GitHub}

GitHub needs your \emph{public} key to verify your identity. Display its contents with:

\begin{lstlisting}[style=bashstyle, caption={Display the public key},escapeinside={(*@}{@*)}]
cat ~/.ssh/id_ed25519.pub
\end{lstlisting}

The output will be a single long line beginning with \texttt{ssh-ed25519}. Select and copy the entire line, including the email address at the end.

Next, navigate to GitHub in your browser and follow these steps:

\begin{enumerate}
    \item Click your profile photo in the upper-right corner and select \textbf{Settings}.
    \item In the left sidebar, click \textbf{SSH and GPG keys}.
    \item Click \textbf{New SSH key}.
    \item Give the key a descriptive title (for example, \textit{Windows laptop}).
    \item Paste the copied public key into the \textbf{Key} field.
    \item Click \textbf{Add SSH key} and confirm with your GitHub password if prompted.
\end{enumerate}

\subsection{Verifying the Connection}

Before using SSH for any repository operations, confirm that GitHub can authenticate your key:

\begin{lstlisting}[style=bashstyle, caption={Test the SSH connection to GitHub},escapeinside={(*@}{@*)}]
ssh -T git@github.com
\end{lstlisting}

On the first connection, SSH will display a fingerprint and ask whether you trust the host. Type \texttt{yes} and press \textbf{Enter}. If authentication succeeds, you will see a message similar to:

\begin{lstlisting}[style=inlineStyle, caption={Expected response from GitHub},escapeinside={(*@}{@*)}]
Hi username! You've successfully authenticated, but GitHub does
not provide shell access.
\end{lstlisting}

If instead you receive a \texttt{Permission denied} error, verify that the agent is running (\texttt{ssh-add -l} should list your key) and that the correct public key has been added to your GitHub account.

\subsection{Cloning a Repository Using SSH}

Once authentication is confirmed, clone repositories using the SSH remote URL:

\begin{lstlisting}[style=bashstyle, caption={Clone a repository over SSH},escapeinside={(*@}{@*)}]
git clone git@github.com:USER/REPOSITORY.git
\end{lstlisting}

Replace \texttt{USER} with the repository owner's GitHub username and \texttt{REPOSITORY} with the repository name.

\begin{warningbox}{Use the SSH Remote URL, Not HTTPS}
The SSH remote URL begins with \texttt{git@github.com:} and ends with \texttt{.git}. It is \emph{not} the same as the HTTPS URL, which begins with \texttt{https://github.com/}. If you clone or push using the HTTPS form after configuring SSH authentication, Git will still prompt you for a GitHub username and password (or personal access token). Always verify the remote URL with \texttt{git remote -v} and update it if necessary:

\begin{lstlisting}[style=bashstyle,escapeinside={(*@}{@*)}]
git remote set-url origin git@github.com:USER/REPOSITORY.git
\end{lstlisting}
\end{warningbox}

\begin{table}[htbp]
\centering
\caption{SSH setup comparison: Linux/macOS vs.\ Windows (Git Bash).}
\label{tab:ssh-comparison}
\begin{tabular}{p{3.5cm} p{5.5cm} p{5.5cm}}
\toprule
\textbf{Step} & \textbf{Linux / macOS} & \textbf{Windows (Git Bash)} \\
\midrule
Open terminal &
Open the system terminal application (Terminal, iTerm2, etc.) &
Right-click Desktop or folder and select \textit{Git Bash Here}, or search Git Bash in the Start menu \\
\addlinespace
Generate key &
\lstinline[style=inlineStyle]{ssh-keygen -t ed25519 -C "email"} (stored in \texttt{\textasciitilde/.ssh/}) &
Same command in Git Bash; default path is \texttt{C:\textbackslash Users\textbackslash Name\textbackslash .ssh\textbackslash} \\
\addlinespace
Start agent &
Agent is typically started automatically by the desktop session; run \lstinline[style=inlineStyle]{eval "$(ssh-agent -s)"} if needed &
Must be started manually each session: \lstinline[style=inlineStyle]{eval "$(ssh-agent -s)"} \\
\addlinespace
Add key to agent &
\lstinline[style=inlineStyle]{ssh-add ~/.ssh/id_ed25519} &
Same command in Git Bash \\
\addlinespace
Display public key &
\lstinline[style=inlineStyle]{cat ~/.ssh/id_ed25519.pub} &
Same command in Git Bash \\
\addlinespace
Verify connection &
\lstinline[style=inlineStyle]{ssh -T git@github.com} &
Same command in Git Bash \\
\addlinespace
Clone with SSH &
\lstinline[style=inlineStyle]{git clone git@github.com:USER/REPO.git} &
Same command in Git Bash \\
\bottomrule
\end{tabular}
\end{table}

% ============================================================
\subsection{Pushing to Multiple Remote Repositories}
\label{subsec:multiple-remotes}
% ============================================================

A single local repository can maintain simultaneous connections to any
number of remote repositories. This is useful when a project must be
mirrored across an organisation account and a personal account (for
example, keeping the \ProjectName{} workshop materials synchronised
between the \texttt{\ProjectName{}} organisation and the personal
\texttt{YourVeryOwnGitHubAccount} GitHub account).

\subsubsection{Concepts}

Git identifies remote repositories by short \emph{names}. The name
\texttt{origin} is the conventional default assigned when a repository
is first cloned. Additional remotes can be registered under any name.
Each remote entry stores a fetch URL and a push URL, both of which can
be set independently.

\begin{infobox}
A remote name is purely local: it exists only in the
\texttt{.git/config} file on your machine and is never pushed to any
server. Two developers cloning the same repository can give the same
remote entirely different names without consequence.
\end{infobox}

\subsubsection{Listing Registered Remotes}

Before adding a remote, confirm what is already registered:

\begin{lstlisting}[style=bashstyle, caption={List all remotes and their URLs},escapeinside={(*@}{@*)}]
git remote -v
\end{lstlisting}

Each remote appears twice, once for fetch and once for push. A freshly
cloned repository typically shows only \texttt{origin}:

\begin{lstlisting}[style=inlineStyle,escapeinside={(*@}{@*)}]
origin  (*@\ProjectGit@*).git (fetch)
origin  (*@\ProjectGit@*).git (push)
\end{lstlisting}

\subsubsection{Adding a Second Remote}

Register an additional remote with \texttt{git remote add}:

\begin{lstlisting}[style=bashstyle, caption={Register a personal mirror remote (Linux / WSL)},escapeinside={(*@}{@*)}]
git remote add YourVeryOwnGitHubAccount \
    (*@\ProjectGit@*).git
\end{lstlisting}

\begin{lstlisting}[style=psstyle, caption={Register a personal mirror remote (PowerShell)},escapeinside={(*@}{@*)}]
git remote add YourVeryOwnGitHubAccount `
    (*@\ProjectGit@*).git
\end{lstlisting}

Verify the result:

\begin{lstlisting}[style=bashstyle,escapeinside={(*@}{@*)}]
git remote -v
\end{lstlisting}

The output should now show four lines:

\begin{lstlisting}[style=inlineStyle,escapeinside={(*@}{@*)}]
YourVeryOwnGitHubAccount  https://github.com/YourVeryOwnGitHubAccount/(*@\ProjectName@*).git (fetch)
YourVeryOwnGitHubAccount  https://github.com/YourVeryOwnGitHubAccount/(*@\ProjectName@*).git (push)
origin           (*@\ProjectGit@*).git (fetch)
origin           (*@\ProjectGit@*).git (push)
\end{lstlisting}

\subsubsection{Pushing to Each Remote}

With two remotes registered, every push must be issued explicitly to
each destination. Git does not push to multiple remotes automatically
unless a push target is configured:

\begin{lstlisting}[style=bashstyle, caption={Push to both remotes explicitly},escapeinside={(*@}{@*)}]
git push origin HEAD
git push YourVeryOwnGitHubAccount HEAD
\end{lstlisting}

\texttt{HEAD} resolves to the current branch, which is the safest form
to use in scripts because it does not depend on knowing the branch name
in advance.

\subsubsection{Checking Whether a Remote Exists Before Pushing}

In automated scripts it is good practice to verify that a remote is
registered before attempting to push to it, so a missing remote produces
a clear diagnostic rather than a fatal error.

\begin{lstlisting}[style=bashstyle, caption={Guard a push with a remote existence check (bash)},escapeinside={(*@}{@*)}]
if git remote get-url YourVeryOwnGitHubAccount > /dev/null 2>&1; then
    git push YourVeryOwnGitHubAccount HEAD
else
    echo "Remote 'YourVeryOwnGitHubAccount' not configured."
    echo "Register it with:"
    echo "  git remote add YourVeryOwnGitHubAccount \
https://github.com/YourVeryOwnGitHubAccount/(*@\ProjectName@*).git"
fi
\end{lstlisting}

\begin{lstlisting}[style=psstyle, caption={Guard a push with a remote existence check (PowerShell / batch)},escapeinside={(*@}{@*)}]
git remote get-url YourVeryOwnGitHubAccount >nul 2>&1
if errorlevel 1 (
    echo Remote 'YourVeryOwnGitHubAccount' not configured.
    echo Register it with:
    echo   git remote add YourVeryOwnGitHubAccount ^
        https://github.com/YourVeryOwnGitHubAccount/(*@\ProjectName@*).git
) else (
    git push YourVeryOwnGitHubAccount HEAD
)
\end{lstlisting}

\subsubsection{Managing Remotes}

Table~\ref{tab:git-remote-commands} summarises the full set of remote
management commands.

\begin{table}[h]
\centering
\caption{Git remote management commands}
\label{tab:git-remote-commands}
\begin{tabular}{lp{8.5cm}}
\toprule
Command & Effect \\
\midrule
\texttt{git remote -v} &
  List all remotes with their fetch and push URLs. \\
\texttt{git remote add <n> <url>} &
  Register a new remote under the given name. \\
\texttt{git remote remove <n>} &
  Delete a remote entry from the local configuration. \\
\texttt{git remote rename <old> <new>} &
  Rename a remote without changing its URL. \\
\texttt{git remote set-url <n> <url>} &
  Update the URL of an existing remote. \\
\texttt{git remote get-url <n>} &
  Print the URL of a named remote; exits with an error if
  the remote does not exist (useful in scripts). \\
\texttt{git remote show <n>} &
  Display detailed information about a remote, including
  tracked branches and push configuration. \\
\bottomrule
\end{tabular}
\end{table}

\subsubsection{Alternative: a Single Remote with Multiple Push URLs}

Git also supports configuring one remote name to push to two URLs
simultaneously, so that a single \texttt{git push origin HEAD} reaches
both destinations. This is convenient but reduces visibility; a failure
on the second URL can be easy to miss in the output.

\begin{lstlisting}[style=bashstyle, caption={Add a second push URL to the origin remote},escapeinside={(*@}{@*)}]
# The first push URL is already set from the clone.
# Add the personal repo as an additional push target:
git remote set-url --add --push origin \
    (*@\ProjectGit@*).git
git remote set-url --add --push origin \
    https://github.com/YourVeryOwnGitHubAccount/(*@\ProjectName@*).git
\end{lstlisting}

\begin{warningbox}
When using \texttt{set-url --add --push}, note that the \emph{first}
call replaces the default push URL rather than adding to it. Both URLs
must be specified explicitly as shown above. Verify the result with
\texttt{git remote -v} before relying on this configuration in
automated scripts.
\end{warningbox}

For the \textsc{Alice} project, the explicit two-remote approach
(separate \texttt{git push} calls per remote) is preferred because
\texttt{gh.bat} checks \texttt{errorlevel} after each push and can
report failures independently. A combined push URL would surface only
the second failure.

% ============================================================
\section{Summary}
\label{sec:git-summary}
% ============================================================

This chapter introduced version control with Git and GitHub, focusing on
the practical steps required to authenticate, clone, and synchronise
repositories.

\begin{itemize}
    \item \textbf{SSH key-based authentication.} GitHub no longer supports
    password authentication for Git operations. This chapter covered the
    complete SSH setup procedure on both Linux and Windows (using Git Bash),
    including key generation with the \texttt{ed25519} algorithm, adding
    the public key to GitHub, and verifying the connection.

    \item \textbf{Cloning repositories via SSH.} Once authentication is
    configured, repositories are cloned using the \texttt{git@github.com:}
    URL form, which avoids credential prompts on every push and pull.

    \item \textbf{Multiple remote repositories.} A single local repository
    can push to more than one remote. This chapter demonstrated how to
    register a second remote and push to both, keeping organisation and
    personal copies in sync.

    \item \textbf{Automated staging, committing, and pushing with
    \texttt{gh.bat}.} The \texttt{gh.bat} script automates the
    add-commit-push cycle, including version stamping and pushing to
    both remotes, reducing repetitive manual steps to a single command.
\end{itemize}
